% ============================================================================
% sections/sec6_conclusion.tex
% ============================================================================
\section{Conclusion and Future Work}
\label{sec:conclusion}

\subsection{Conclusion}

A layered, empirically honest defense architecture against retrieval
poisoning was designed, implemented, and evaluated against real attack
constructions drawn from six cited papers. No single ring is sufficient
alone --- each has a real, measured limitation, stated in this report
rather than smoothed over --- but the combination of four rings plus a
persistent, cross-query trust store closes the measured gaps within this
simulation's scope, reaching $100.0\pm0.0\%$ accuracy and $0.0\pm0.0\%$
overall attack success across 1{,}600 evaluation queries
(\Cref{tab:main-comparison}).

Equally significant for a project of this kind is its debugging
discipline. Eight distinct, real, measured bugs were found across the
project's development --- from a self-referential filter threshold to a
consensus mechanism that silently did nothing at all
(\Cref{tab:debugging}) --- each fixed by direct measurement rather than
by adjusting a parameter until a target number appeared. That
verification trail is itself evidence of sound methodology and is
presented here alongside the final numbers rather than left out of the
report.

The most genuinely novel contributions of this work are, in order of how
directly they were earned through this project's own measurement rather
than inherited from the surveyed literature:

\begin{enumerate}
\item Dual-signal risk routing (embedding cohesion together with
  answer-vote contention) --- discovered empirically during this project,
  not inherited directly from any one cited paper.
\item Group-wise counterfactual consensus --- a genuine generalization of
  RAGuard's leave-one-out to colluding cliques, with an honestly reported
  single-query ceiling (\Cref{sec:discussion}), reached only after a real
  implementation bug was found and fixed (\Cref{sec:debugging}, bug \#7).
\item A persistent, cross-query trust store --- closes the specific
  single-query information ceiling that \Cref{sec:discussion} identifies
  GWCC alone cannot close.
\item Risk-based compute escalation --- a fast path for the common case,
  with deep validation reserved for queries that genuinely trigger a risk
  signal, rather than a fixed multi-pass cost on every query.
\item A unified, multi-threat evaluation harness --- six attack regimes
  against six systems, with multi-seed confidence intervals and a
  per-ring ablation, in one framework built specifically so its
  components cannot silently drift apart from one another
  (\Cref{sec:multiseed}).
\end{enumerate}

\subsection{Future Work}
\label{sec:future-work}

\begin{enumerate}
\item Run the TF-IDF-vs-LSA embedding comparison already present in the
  codebase (\texttt{run\_\allowbreak embedding\_\allowbreak
  comparison.py}) to test whether Ring~1's vocabulary-novelty advantage
  survives a richer embedding space.
\item Replace \texttt{doc.answer} with a real, small LLM call, to
  validate end-to-end the one assumption every ring in this pipeline
  currently depends on --- that a document's supported answer is reliably
  extractable by a language model, not only by a ground-truth label.
\item Test against adaptive attackers specifically optimized to evade
  Ring~0's repetition-ratio threshold and Ring~2's contention threshold,
  rather than the realistic-but-non-adaptive attacks evaluated here.
\item Explicitly measure the detection-versus-runtime trade-off of
  Ring~3's \texttt{MAX\_PAIR\_SAMPLES} cap, rather than leaving it fixed
  at its current constant.
\item Validate the whole pipeline on a naturally authored corpus (for
  example a BEIR, Natural Questions, or HotpotQA subset), to check
  whether the shared-generator finding in \Cref{sec:shared-generator} ---
  that overly uniform clean text can create unintended structural
  fingerprints --- generalizes beyond this project's synthetic,
  template-generated text.
\item Package the benchmark harness itself as an open-source
  red-teaming tool, so the specific question this project's own
  \Cref{sec:discussion} raises --- does a given RAG defense's risk signal
  measure more than one attack axis --- can be tested against other,
  independently built RAG pipelines, not only against the baselines
  re-implemented here.
\end{enumerate}
